%%%%%%%%%%%%%%%%%%%%%%%%%%%%%%%%%%%%%%%%%%%%%%%%%%%%%%%%%%%%
% 10.15 GAME THEORY
% The Composition of Advanced Mathematics
%%%%%%%%%%%%%%%%%%%%%%%%%%%%%%%%%%%%%%%%%%%%%%%%%%%%%%%%%%%%

\section{Game Theory}
\label{sec:game-theory}

\prerequisite{
Decision theory, probability, optimization fundamentals, linear
programming, expected utility, matrices, calculus, and basic economics.
}

\learningobjectives{
By the end of this section, the reader should be able to:
\begin{itemize}
    \item define a strategic game;
    \item identify players, strategies, actions, and payoffs;
    \item distinguish cooperative and noncooperative games;
    \item distinguish simultaneous and sequential games;
    \item represent normal-form games using payoff matrices;
    \item identify dominant and dominated strategies;
    \item apply iterated elimination of strictly dominated strategies;
    \item define best responses;
    \item define Nash equilibrium;
    \item identify pure-strategy Nash equilibria;
    \item understand coordination and anti-coordination games;
    \item analyze the Prisoner's Dilemma;
    \item analyze Battle of the Sexes;
    \item analyze Matching Pennies;
    \item define mixed strategies;
    \item compute expected payoffs under mixed strategies;
    \item solve simple mixed-strategy Nash equilibria;
    \item define zero-sum games;
    \item distinguish maximin and minimax values;
    \item state the minimax theorem;
    \item formulate finite zero-sum games as linear programs;
    \item understand security strategies;
    \item distinguish constant-sum and general-sum games;
    \item represent sequential games in extensive form;
    \item identify decision nodes and information sets;
    \item use backward induction;
    \item define subgames;
    \item define subgame-perfect equilibrium;
    \item understand credible and noncredible threats;
    \item recognize commitment and first-mover effects;
    \item analyze simple entry-deterrence games;
    \item understand repeated games conceptually;
    \item distinguish finite and infinite repetition;
    \item understand discounting in repeated games;
    \item recognize trigger strategies;
    \item understand the Folk Theorem conceptually;
    \item recognize Bayesian games;
    \item define types and beliefs;
    \item understand Bayesian Nash equilibrium conceptually;
    \item understand signaling and screening conceptually;
    \item recognize auctions as strategic games;
    \item understand first-price and second-price sealed-bid auctions;
    \item recognize truthful bidding in second-price auctions;
    \item understand mechanism-design concepts at an introductory level;
    \item define Pareto efficiency in strategic settings;
    \item distinguish Nash equilibrium from social optimality;
    \item understand price of anarchy conceptually;
    \item recognize congestion games;
    \item define potential games conceptually;
    \item understand evolutionary game theory conceptually;
    \item recognize evolutionary stable strategies;
    \item connect game theory with economics, operations research,
          networks, markets, auctions, politics, biology, and distributed
          systems.
\end{itemize}
}

%%%%%%%%%%%%%%%%%%%%%%%%%%%%%%%%%%%%%%%%%%%%%%%%%%%%%%%%%%%%
\subsection{What Is Game Theory?}
\label{subsec:what-is-game-theory}
%%%%%%%%%%%%%%%%%%%%%%%%%%%%%%%%%%%%%%%%%%%%%%%%%%%%%%%%%%%%

Game theory studies decision problems in which the outcome for each
decision maker depends not only on their own choice, but also on the
choices of other decision makers.

The participants are called

\[
\boxed{
\text{players}.
}
\]

A strategic interaction contains:

\[
\boxed{
\text{players}
+
\text{strategies}
+
\text{payoffs}.
}
\]

%%%%%%%%%%%%%%%%%%%%%%%%%%%%%%%%%%%%%%%%%%%%%%%%%%%%%%%%%%%%
\subsection{Strategic Interaction}
\label{subsec:strategic-interaction}
%%%%%%%%%%%%%%%%%%%%%%%%%%%%%%%%%%%%%%%%%%%%%%%%%%%%%%%%%%%%

In ordinary decision theory, uncertainty may come from nature.

In game theory, uncertainty often comes from the deliberate choices of
other agents.

Each player therefore asks:

\[
\boxed{
\text{What should I do given what the others may do?}
}
\]

Each player's optimal choice may depend on the choices of others.

%%%%%%%%%%%%%%%%%%%%%%%%%%%%%%%%%%%%%%%%%%%%%%%%%%%%%%%%%%%%
\subsection{Players}
\label{subsec:game-theory-players}
%%%%%%%%%%%%%%%%%%%%%%%%%%%%%%%%%%%%%%%%%%%%%%%%%%%%%%%%%%%%

Let

\[
\boxed{
N
=
\{1,2,\ldots,n\}
}
\]

be the set of players.

Examples include:

\[
\boxed{
\text{firms},
}
\]

\[
\boxed{
\text{countries},
}
\]

\[
\boxed{
\text{drivers},
}
\]

\[
\boxed{
\text{bidders},
}
\]

or

\[
\boxed{
\text{computer agents}.
}
\]

%%%%%%%%%%%%%%%%%%%%%%%%%%%%%%%%%%%%%%%%%%%%%%%%%%%%%%%%%%%%
\subsection{Strategies}
\label{subsec:game-theory-strategies}
%%%%%%%%%%%%%%%%%%%%%%%%%%%%%%%%%%%%%%%%%%%%%%%%%%%%%%%%%%%%

Player \(i\) has a strategy set

\[
\boxed{
S_i.
}
\]

A strategy

\[
s_i\in S_i
\]

specifies the player's behavior.

The full strategy profile is

\[
\boxed{
s
=
(s_1,s_2,\ldots,s_n).
}
\]

The strategies of all players except \(i\) are denoted

\[
\boxed{
s_{-i}.
}
\]

%%%%%%%%%%%%%%%%%%%%%%%%%%%%%%%%%%%%%%%%%%%%%%%%%%%%%%%%%%%%
\subsection{Payoff Functions}
\label{subsec:game-theory-payoffs}
%%%%%%%%%%%%%%%%%%%%%%%%%%%%%%%%%%%%%%%%%%%%%%%%%%%%%%%%%%%%

Player \(i\)'s payoff is

\[
\boxed{
u_i(s_1,\ldots,s_n).
}
\]

Each player typically seeks to maximize their own payoff.

Thus the strategic problem faced by player \(i\) is

\[
\boxed{
\max_{s_i\in S_i}
u_i(s_i,s_{-i}).
}
\]

%%%%%%%%%%%%%%%%%%%%%%%%%%%%%%%%%%%%%%%%%%%%%%%%%%%%%%%%%%%%
\subsection{Normal-Form Games}
\label{subsec:normal-form-games}
%%%%%%%%%%%%%%%%%%%%%%%%%%%%%%%%%%%%%%%%%%%%%%%%%%%%%%%%%%%%

A normal-form game is specified by

\[
\boxed{
G
=
\left(
N,
(S_i)_{i\in N},
(u_i)_{i\in N}
\right).
}
\]

For two players with finitely many actions, the game can be represented by
a payoff matrix.

%%%%%%%%%%%%%%%%%%%%%%%%%%%%%%%%%%%%%%%%%%%%%%%%%%%%%%%%%%%%
\subsection{Payoff Matrix}
\label{subsec:game-payoff-matrix}
%%%%%%%%%%%%%%%%%%%%%%%%%%%%%%%%%%%%%%%%%%%%%%%%%%%%%%%%%%%%

Consider two players.

Player \(1\) chooses a row.

Player \(2\) chooses a column.

A payoff matrix may be written

\[
\begin{array}{c|cc}
&
L
&
R
\\
\hline
U
&
(a_{11},b_{11})
&
(a_{12},b_{12})
\\
D
&
(a_{21},b_{21})
&
(a_{22},b_{22})
\end{array}
\]

The first number in each cell is Player \(1\)'s payoff.

The second is Player \(2\)'s payoff.

%%%%%%%%%%%%%%%%%%%%%%%%%%%%%%%%%%%%%%%%%%%%%%%%%%%%%%%%%%%%
\subsection{Cooperative and Noncooperative Games}
\label{subsec:cooperative-noncooperative-games}
%%%%%%%%%%%%%%%%%%%%%%%%%%%%%%%%%%%%%%%%%%%%%%%%%%%%%%%%%%%%

Noncooperative game theory studies strategic choices made by individual
players under specified rules.

Cooperative game theory studies what groups of players can achieve through
coalitions and how coalition value should be allocated.

This section focuses primarily on noncooperative game theory.

%%%%%%%%%%%%%%%%%%%%%%%%%%%%%%%%%%%%%%%%%%%%%%%%%%%%%%%%%%%%
\subsection{Simultaneous and Sequential Games}
\label{subsec:simultaneous-sequential-games}
%%%%%%%%%%%%%%%%%%%%%%%%%%%%%%%%%%%%%%%%%%%%%%%%%%%%%%%%%%%%

In a simultaneous game, players choose without observing the current choice
of the other players.

In a sequential game, later players may observe earlier moves before
choosing.

Thus:

\[
\boxed{
\text{simultaneous games}
\rightarrow
\text{normal form},
}
\]

while

\[
\boxed{
\text{sequential games}
\rightarrow
\text{extensive form}.
}
\]

%%%%%%%%%%%%%%%%%%%%%%%%%%%%%%%%%%%%%%%%%%%%%%%%%%%%%%%%%%%%
\subsection{Dominant Strategies}
\label{subsec:dominant-strategies}
%%%%%%%%%%%%%%%%%%%%%%%%%%%%%%%%%%%%%%%%%%%%%%%%%%%%%%%%%%%%

A strategy \(s_i^*\) is strictly dominant for player \(i\) if

\[
\boxed{
u_i(s_i^*,s_{-i})
>
u_i(s_i,s_{-i})
}
\]

for every alternative

\[
s_i\neq s_i^*
\]

and every strategy profile \(s_{-i}\) of the other players.

A dominant strategy is optimal regardless of what the other players do.

%%%%%%%%%%%%%%%%%%%%%%%%%%%%%%%%%%%%%%%%%%%%%%%%%%%%%%%%%%%%
\subsection{Weak Dominance}
\label{subsec:weak-dominance}
%%%%%%%%%%%%%%%%%%%%%%%%%%%%%%%%%%%%%%%%%%%%%%%%%%%%%%%%%%%%

Strategy \(s_i^*\) weakly dominates \(s_i\) if

\[
\boxed{
u_i(s_i^*,s_{-i})
\geq
u_i(s_i,s_{-i})
}
\]

for every \(s_{-i}\), with strict inequality for at least one opponent
strategy.

Weak dominance must be handled more carefully than strict dominance.

%%%%%%%%%%%%%%%%%%%%%%%%%%%%%%%%%%%%%%%%%%%%%%%%%%%%%%%%%%%%
\subsection{Dominated Strategies}
\label{subsec:dominated-strategies}
%%%%%%%%%%%%%%%%%%%%%%%%%%%%%%%%%%%%%%%%%%%%%%%%%%%%%%%%%%%%

A strategy is strictly dominated if another strategy produces a strictly
larger payoff for every possible strategy of the opponents.

A rational player should never need a strictly dominated strategy when
maximizing payoff.

This motivates elimination procedures.

%%%%%%%%%%%%%%%%%%%%%%%%%%%%%%%%%%%%%%%%%%%%%%%%%%%%%%%%%%%%
\subsection{Iterated Elimination of Strictly Dominated Strategies}
\label{subsec:iterated-elimination}
%%%%%%%%%%%%%%%%%%%%%%%%%%%%%%%%%%%%%%%%%%%%%%%%%%%%%%%%%%%%

The procedure is:

\begin{enumerate}
    \item identify a strictly dominated strategy;
    \item remove it;
    \item reconsider dominance in the reduced game;
    \item continue until no additional strictly dominated strategies
          remain.
\end{enumerate}

This process can simplify large payoff matrices.

%%%%%%%%%%%%%%%%%%%%%%%%%%%%%%%%%%%%%%%%%%%%%%%%%%%%%%%%%%%%
\subsection{Worked Example: Dominant Strategy}
\label{subsec:worked-dominant-strategy}
%%%%%%%%%%%%%%%%%%%%%%%%%%%%%%%%%%%%%%%%%%%%%%%%%%%%%%%%%%%%

\begin{workedexamplebox}{Finding a Dominant Strategy}

Consider:

\[
\begin{array}{c|cc}
&
L
&
R
\\
\hline
U
&
(5,2)
&
(4,4)
\\
D
&
(2,3)
&
(1,1)
\end{array}
\]

For Player \(1\):

if Player \(2\) chooses \(L\),

\[
5>2.
\]

If Player \(2\) chooses \(R\),

\[
4>1.
\]

Thus \(U\) strictly dominates \(D\).

Therefore,

\begin{answerbox}
\[
\boxed{
U
\text{ is a strictly dominant strategy for Player }1.
}
\]
\end{answerbox}

\end{workedexamplebox}

%%%%%%%%%%%%%%%%%%%%%%%%%%%%%%%%%%%%%%%%%%%%%%%%%%%%%%%%%%%%
\subsection{Best Responses}
\label{subsec:best-responses}
%%%%%%%%%%%%%%%%%%%%%%%%%%%%%%%%%%%%%%%%%%%%%%%%%%%%%%%%%%%%

Player \(i\)'s best-response correspondence is

\[
\boxed{
BR_i(s_{-i})
=
\operatorname*{arg\,max}_{s_i\in S_i}
u_i(s_i,s_{-i}).
}
\]

A best response is optimal given the opponents' strategies.

Unlike a dominant strategy, it may depend on what the opponents choose.

%%%%%%%%%%%%%%%%%%%%%%%%%%%%%%%%%%%%%%%%%%%%%%%%%%%%%%%%%%%%
\subsection{Nash Equilibrium}
\label{subsec:nash-equilibrium}
%%%%%%%%%%%%%%%%%%%%%%%%%%%%%%%%%%%%%%%%%%%%%%%%%%%%%%%%%%%%

\begin{definitionbox}{Nash Equilibrium}
A strategy profile

\[
s^*
=
(s_1^*,\ldots,s_n^*)
\]

is a Nash equilibrium if every player's strategy is a best response to the
strategies of all other players.
\end{definitionbox}

Thus,

\[
\boxed{
s_i^*
\in
BR_i(s_{-i}^*)
}
\]

for every player \(i\).

%%%%%%%%%%%%%%%%%%%%%%%%%%%%%%%%%%%%%%%%%%%%%%%%%%%%%%%%%%%%
\subsection{Nash Equilibrium Inequality}
\label{subsec:nash-equilibrium-inequality}
%%%%%%%%%%%%%%%%%%%%%%%%%%%%%%%%%%%%%%%%%%%%%%%%%%%%%%%%%%%%

Equivalently,

\[
\boxed{
u_i(s_i^*,s_{-i}^*)
\geq
u_i(s_i,s_{-i}^*)
}
\]

for every

\[
s_i\in S_i
\]

and every player \(i\).

No single player can improve their payoff through unilateral deviation.

%%%%%%%%%%%%%%%%%%%%%%%%%%%%%%%%%%%%%%%%%%%%%%%%%%%%%%%%%%%%
\subsection{Nash Equilibrium Is Not Joint Optimality}
\label{subsec:nash-not-joint-optimality}
%%%%%%%%%%%%%%%%%%%%%%%%%%%%%%%%%%%%%%%%%%%%%%%%%%%%%%%%%%%%

A Nash equilibrium is stable against unilateral deviations.

It need not maximize:

\[
\boxed{
\sum_i
u_i.
}
\]

It also need not be Pareto efficient.

This distinction is fundamental.

%%%%%%%%%%%%%%%%%%%%%%%%%%%%%%%%%%%%%%%%%%%%%%%%%%%%%%%%%%%%
\subsection{Finding Pure-Strategy Nash Equilibria}
\label{subsec:finding-pure-nash}
%%%%%%%%%%%%%%%%%%%%%%%%%%%%%%%%%%%%%%%%%%%%%%%%%%%%%%%%%%%%

For a finite two-player matrix game:

\begin{enumerate}
    \item identify Player \(1\)'s best response in each column;
    \item identify Player \(2\)'s best response in each row;
    \item locate cells that are best responses for both players.
\end{enumerate}

Those cells are pure-strategy Nash equilibria.

%%%%%%%%%%%%%%%%%%%%%%%%%%%%%%%%%%%%%%%%%%%%%%%%%%%%%%%%%%%%
\subsection{Worked Example: Pure Nash Equilibrium}
\label{subsec:worked-pure-nash}
%%%%%%%%%%%%%%%%%%%%%%%%%%%%%%%%%%%%%%%%%%%%%%%%%%%%%%%%%%%%

\begin{workedexamplebox}{Best-Response Intersection}

Consider

\[
\begin{array}{c|cc}
&
L
&
R
\\
\hline
U
&
(3,3)
&
(0,2)
\\
D
&
(2,0)
&
(1,1)
\end{array}
\]

If Player \(2\) chooses \(L\), Player \(1\) prefers \(U\):

\[
3>2.
\]

If Player \(2\) chooses \(R\), Player \(1\) prefers \(D\):

\[
1>0.
\]

If Player \(1\) chooses \(U\), Player \(2\) prefers \(L\):

\[
3>2.
\]

If Player \(1\) chooses \(D\), Player \(2\) prefers \(R\):

\[
1>0.
\]

Thus there are two mutual best-response cells:

\begin{answerbox}
\[
\boxed{
(U,L)
\quad\text{and}\quad
(D,R).
}
\]
\end{answerbox}

\end{workedexamplebox}

%%%%%%%%%%%%%%%%%%%%%%%%%%%%%%%%%%%%%%%%%%%%%%%%%%%%%%%%%%%%
\subsection{Prisoner's Dilemma}
\label{subsec:prisoners-dilemma}
%%%%%%%%%%%%%%%%%%%%%%%%%%%%%%%%%%%%%%%%%%%%%%%%%%%%%%%%%%%%

A canonical Prisoner's Dilemma has payoff structure

\[
\begin{array}{c|cc}
&
C
&
D
\\
\hline
C
&
(R,R)
&
(S,T)
\\
D
&
(T,S)
&
(P,P)
\end{array}
\]

with

\[
\boxed{
T>R>P>S.
}
\]

Here:

\[
C
=
\text{cooperate},
\]

and

\[
D
=
\text{defect}.
}
\]

%%%%%%%%%%%%%%%%%%%%%%%%%%%%%%%%%%%%%%%%%%%%%%%%%%%%%%%%%%%%
\subsection{Prisoner's Dilemma Equilibrium}
\label{subsec:prisoners-dilemma-equilibrium}
%%%%%%%%%%%%%%%%%%%%%%%%%%%%%%%%%%%%%%%%%%%%%%%%%%%%%%%%%%%%

Defection is strictly dominant when

\[
T>R
\]

and

\[
P>S.
\]

Therefore the unique Nash equilibrium is

\[
\boxed{
(D,D).
}
\]

Yet both players prefer

\[
(C,C)
\]

to

\[
(D,D)
\]

because

\[
R>P.
\]

This demonstrates that Nash equilibrium need not be Pareto efficient.

%%%%%%%%%%%%%%%%%%%%%%%%%%%%%%%%%%%%%%%%%%%%%%%%%%%%%%%%%%%%
\subsection{Coordination Games}
\label{subsec:coordination-games}
%%%%%%%%%%%%%%%%%%%%%%%%%%%%%%%%%%%%%%%%%%%%%%%%%%%%%%%%%%%%

A coordination game rewards players for choosing compatible actions.

For example,

\[
\begin{array}{c|cc}
&
L
&
R
\\
\hline
L
&
(2,2)
&
(0,0)
\\
R
&
(0,0)
&
(1,1)
\end{array}
\]

has two pure Nash equilibria:

\[
\boxed{
(L,L)
}
\]

and

\[
\boxed{
(R,R).
}
\]

The strategic difficulty is equilibrium selection.

%%%%%%%%%%%%%%%%%%%%%%%%%%%%%%%%%%%%%%%%%%%%%%%%%%%%%%%%%%%%
\subsection{Battle of the Sexes}
\label{subsec:battle-of-the-sexes}
%%%%%%%%%%%%%%%%%%%%%%%%%%%%%%%%%%%%%%%%%%%%%%%%%%%%%%%%%%%%

Battle of the Sexes is a coordination game in which players prefer to
coordinate but disagree about which coordinated outcome is best.

A typical payoff matrix is

\[
\begin{array}{c|cc}
&
O
&
F
\\
\hline
O
&
(2,1)
&
(0,0)
\\
F
&
(0,0)
&
(1,2)
\end{array}
\]

with pure equilibria

\[
\boxed{
(O,O)
}
\]

and

\[
\boxed{
(F,F).
}
\]

%%%%%%%%%%%%%%%%%%%%%%%%%%%%%%%%%%%%%%%%%%%%%%%%%%%%%%%%%%%%
\subsection{Anti-Coordination Games}
\label{subsec:anti-coordination-games}
%%%%%%%%%%%%%%%%%%%%%%%%%%%%%%%%%%%%%%%%%%%%%%%%%%%%%%%%%%%%

In anti-coordination games, players prefer different actions.

Examples include:

\[
\boxed{
\text{choosing opposite sides of a congestion problem}
}
\]

or

\[
\boxed{
\text{specializing in different roles}.
}
\]

The strategic incentive favors separation rather than matching.

%%%%%%%%%%%%%%%%%%%%%%%%%%%%%%%%%%%%%%%%%%%%%%%%%%%%%%%%%%%%
\subsection{Matching Pennies}
\label{subsec:matching-pennies}
%%%%%%%%%%%%%%%%%%%%%%%%%%%%%%%%%%%%%%%%%%%%%%%%%%%%%%%%%%%%

Matching Pennies is a two-player zero-sum game:

\[
\begin{array}{c|cc}
&
H
&
T
\\
\hline
H
&
(1,-1)
&
(-1,1)
\\
T
&
(-1,1)
&
(1,-1)
\end{array}
\]

No cell is a mutual best response.

Thus there is no pure-strategy Nash equilibrium.

A mixed strategy is required.

%%%%%%%%%%%%%%%%%%%%%%%%%%%%%%%%%%%%%%%%%%%%%%%%%%%%%%%%%%%%
\subsection{Mixed Strategies}
\label{subsec:mixed-strategies}
%%%%%%%%%%%%%%%%%%%%%%%%%%%%%%%%%%%%%%%%%%%%%%%%%%%%%%%%%%%%

A mixed strategy assigns probabilities to pure strategies.

If Player \(1\) has pure strategies

\[
s_1,\ldots,s_k,
\]

a mixed strategy is

\[
\boxed{
p
=
(p_1,\ldots,p_k),
}
\]

where

\[
p_i\geq0
\]

and

\[
\boxed{
\sum_i
p_i=1.
}
\]

%%%%%%%%%%%%%%%%%%%%%%%%%%%%%%%%%%%%%%%%%%%%%%%%%%%%%%%%%%%%
\subsection{Expected Payoff Under Mixed Strategies}
\label{subsec:mixed-strategy-expected-payoff}
%%%%%%%%%%%%%%%%%%%%%%%%%%%%%%%%%%%%%%%%%%%%%%%%%%%%%%%%%%%%

For a two-player game, let Player \(1\)'s payoff matrix be \(A\).

If Player \(1\) uses probability vector \(p\) and Player \(2\) uses \(q\),
then Player \(1\)'s expected payoff is

\[
\boxed{
p^TAq.
}
\]

For a zero-sum game, Player \(2\)'s payoff is

\[
\boxed{
-p^TAq.
}
\]

%%%%%%%%%%%%%%%%%%%%%%%%%%%%%%%%%%%%%%%%%%%%%%%%%%%%%%%%%%%%
\subsection{Indifference Principle}
\label{subsec:mixed-strategy-indifference}
%%%%%%%%%%%%%%%%%%%%%%%%%%%%%%%%%%%%%%%%%%%%%%%%%%%%%%%%%%%%

In a mixed-strategy equilibrium, every pure strategy used with positive
probability must give the same expected payoff against the opponent's
equilibrium mixture.

Otherwise the player would shift probability toward a strictly better
pure strategy.

This principle makes many \(2\times2\) mixed equilibria easy to compute.

%%%%%%%%%%%%%%%%%%%%%%%%%%%%%%%%%%%%%%%%%%%%%%%%%%%%%%%%%%%%
\subsection{Worked Example: Matching Pennies Mixed Equilibrium}
\label{subsec:worked-matching-pennies}
%%%%%%%%%%%%%%%%%%%%%%%%%%%%%%%%%%%%%%%%%%%%%%%%%%%%%%%%%%%%

\begin{workedexamplebox}{Mixed Nash Equilibrium}

Suppose Player \(2\) chooses \(H\) with probability \(q\).

If Player \(1\) chooses \(H\), expected payoff is

\[
q(1)+(1-q)(-1)
=
2q-1.
\]

If Player \(1\) chooses \(T\), expected payoff is

\[
q(-1)+(1-q)(1)
=
1-2q.
\]

Player \(1\) is indifferent when

\[
2q-1
=
1-2q.
\]

Thus,

\[
4q=2,
\]

so

\[
q=\frac12.
\]

By symmetry,

\[
p=\frac12.
\]

Therefore,

\begin{answerbox}
\[
\boxed{
p^*
=
q^*
=
\frac12.
}
\]
\end{answerbox}

Each player randomizes equally between heads and tails.

\end{workedexamplebox}

%%%%%%%%%%%%%%%%%%%%%%%%%%%%%%%%%%%%%%%%%%%%%%%%%%%%%%%%%%%%
\subsection{Mixed-Strategy Nash Equilibrium}
\label{subsec:mixed-nash-equilibrium}
%%%%%%%%%%%%%%%%%%%%%%%%%%%%%%%%%%%%%%%%%%%%%%%%%%%%%%%%%%%%

A mixed-strategy profile is a Nash equilibrium if no player can increase
expected payoff by changing their probability distribution unilaterally.

Nash's theorem guarantees that every finite game has at least one Nash
equilibrium in mixed strategies.

%%%%%%%%%%%%%%%%%%%%%%%%%%%%%%%%%%%%%%%%%%%%%%%%%%%%%%%%%%%%
\subsection{Support of a Mixed Strategy}
\label{subsec:mixed-strategy-support}
%%%%%%%%%%%%%%%%%%%%%%%%%%%%%%%%%%%%%%%%%%%%%%%%%%%%%%%%%%%%

The support of a mixed strategy is the set of pure strategies assigned
positive probability.

At equilibrium, pure strategies in the support must be best responses to
the opponents' equilibrium strategy.

Pure strategies outside the support cannot yield strictly larger payoff.

%%%%%%%%%%%%%%%%%%%%%%%%%%%%%%%%%%%%%%%%%%%%%%%%%%%%%%%%%%%%
\subsection{General \(2\times2\) Mixed Equilibrium}
\label{subsec:general-two-by-two-mixed}
%%%%%%%%%%%%%%%%%%%%%%%%%%%%%%%%%%%%%%%%%%%%%%%%%%%%%%%%%%%%

Suppose Player \(1\)'s payoff matrix is

\[
A
=
\begin{pmatrix}
a & b\\
c & d
\end{pmatrix}.
\]

If Player \(2\) chooses the first column with probability \(q\), Player
\(1\)'s expected payoff from row \(1\) is

\[
aq+b(1-q),
\]

while row \(2\) yields

\[
cq+d(1-q).
\]

An interior equilibrium requires

\[
\boxed{
aq+b(1-q)
=
cq+d(1-q).
}
\]

Solve this equation for \(q\).

A corresponding indifference equation determines Player \(1\)'s mixing
probability.

%%%%%%%%%%%%%%%%%%%%%%%%%%%%%%%%%%%%%%%%%%%%%%%%%%%%%%%%%%%%
\subsection{Zero-Sum Games}
\label{subsec:zero-sum-games}
%%%%%%%%%%%%%%%%%%%%%%%%%%%%%%%%%%%%%%%%%%%%%%%%%%%%%%%%%%%%

A two-player zero-sum game satisfies

\[
\boxed{
u_1(s_1,s_2)
+
u_2(s_1,s_2)
=
0.
}
\]

One player's gain is exactly the other player's loss.

Therefore it is sufficient to specify one payoff matrix \(A\).

%%%%%%%%%%%%%%%%%%%%%%%%%%%%%%%%%%%%%%%%%%%%%%%%%%%%%%%%%%%%
\subsection{Security Level}
\label{subsec:security-level}
%%%%%%%%%%%%%%%%%%%%%%%%%%%%%%%%%%%%%%%%%%%%%%%%%%%%%%%%%%%%

Suppose Player \(1\) chooses a pure row \(i\).

The worst payoff Player \(2\) can force is

\[
\min_j
a_{ij}.
\]

Player \(1\)'s pure-strategy security level is therefore

\[
\boxed{
\max_i
\min_j
a_{ij}.
}
\]

This is the maximin value.

%%%%%%%%%%%%%%%%%%%%%%%%%%%%%%%%%%%%%%%%%%%%%%%%%%%%%%%%%%%%
\subsection{Minimax Value}
\label{subsec:minimax-value-zero-sum}
%%%%%%%%%%%%%%%%%%%%%%%%%%%%%%%%%%%%%%%%%%%%%%%%%%%%%%%%%%%%

Player \(2\) seeks to minimize Player \(1\)'s payoff.

For column \(j\), Player \(1\) could obtain at most

\[
\max_i
a_{ij}.
\]

Thus Player \(2\)'s pure minimax value is

\[
\boxed{
\min_j
\max_i
a_{ij}.
}
\]

Always,

\[
\boxed{
\max_i\min_j a_{ij}
\leq
\min_j\max_i a_{ij}.
}
\]

%%%%%%%%%%%%%%%%%%%%%%%%%%%%%%%%%%%%%%%%%%%%%%%%%%%%%%%%%%%%
\subsection{Saddle Point}
\label{subsec:game-saddle-point}
%%%%%%%%%%%%%%%%%%%%%%%%%%%%%%%%%%%%%%%%%%%%%%%%%%%%%%%%%%%%

If

\[
\boxed{
\max_i\min_j a_{ij}
=
\min_j\max_i a_{ij},
}
\]

the corresponding entry is a pure-strategy saddle point.

Then the zero-sum game has a pure-strategy equilibrium.

If the quantities differ, randomization can close the gap.

%%%%%%%%%%%%%%%%%%%%%%%%%%%%%%%%%%%%%%%%%%%%%%%%%%%%%%%%%%%%
\subsection{Von Neumann Minimax Theorem}
\label{subsec:minimax-theorem}
%%%%%%%%%%%%%%%%%%%%%%%%%%%%%%%%%%%%%%%%%%%%%%%%%%%%%%%%%%%%

\begin{theorem}[Minimax Theorem]
For every finite two-player zero-sum game with payoff matrix \(A\),

\[
\boxed{
\max_{p}
\min_{q}
p^TAq
=
\min_{q}
\max_{p}
p^TAq.
}
\]

The common quantity is the value of the game.
\end{theorem}

Mixed strategies therefore guarantee matching maximin and minimax values.

%%%%%%%%%%%%%%%%%%%%%%%%%%%%%%%%%%%%%%%%%%%%%%%%%%%%%%%%%%%%
\subsection{Value of a Zero-Sum Game}
\label{subsec:value-zero-sum-game}
%%%%%%%%%%%%%%%%%%%%%%%%%%%%%%%%%%%%%%%%%%%%%%%%%%%%%%%%%%%%

The value \(v\) is the expected equilibrium payoff to Player \(1\).

Thus Player \(1\) can guarantee at least

\[
\boxed{
v,
}
\]

while Player \(2\) can hold Player \(1\) to at most

\[
\boxed{
v.
}
\]

At equilibrium, both guarantees coincide.

%%%%%%%%%%%%%%%%%%%%%%%%%%%%%%%%%%%%%%%%%%%%%%%%%%%%%%%%%%%%
\subsection{Linear Programming Formulation for Player 1}
\label{subsec:zero-sum-player-one-lp}
%%%%%%%%%%%%%%%%%%%%%%%%%%%%%%%%%%%%%%%%%%%%%%%%%%%%%%%%%%%%

Suppose Player \(1\) chooses mixed strategy \(p\).

To guarantee value \(v\),

\[
\boxed{
\sum_i
p_i a_{ij}
\geq
v
}
\]

for every column \(j\).

Therefore Player \(1\) solves

\[
\boxed{
\begin{aligned}
\max_{p,v}
\quad&
v
\\
\text{subject to}
\quad&
\sum_i
p_i a_{ij}
\geq
v,
\qquad
\text{for every }j,
\\
&
\sum_i
p_i=1,
\\
&
p_i\geq0.
\end{aligned}
}
\]

%%%%%%%%%%%%%%%%%%%%%%%%%%%%%%%%%%%%%%%%%%%%%%%%%%%%%%%%%%%%
\subsection{Linear Programming Formulation for Player 2}
\label{subsec:zero-sum-player-two-lp}
%%%%%%%%%%%%%%%%%%%%%%%%%%%%%%%%%%%%%%%%%%%%%%%%%%%%%%%%%%%%

Player \(2\) chooses probabilities \(q_j\) to minimize an upper bound \(w\)
on Player \(1\)'s payoff:

\[
\boxed{
\begin{aligned}
\min_{q,w}
\quad&
w
\\
\text{subject to}
\quad&
\sum_j
a_{ij}q_j
\leq
w,
\qquad
\text{for every }i,
\\
&
\sum_j
q_j=1,
\\
&
q_j\geq0.
\end{aligned}
}
\]

The two LPs are dual in an appropriate transformed formulation.

%%%%%%%%%%%%%%%%%%%%%%%%%%%%%%%%%%%%%%%%%%%%%%%%%%%%%%%%%%%%
\subsection{Game Theory and Linear Programming Duality}
\label{subsec:game-theory-lp-duality}
%%%%%%%%%%%%%%%%%%%%%%%%%%%%%%%%%%%%%%%%%%%%%%%%%%%%%%%%%%%%

The minimax theorem is closely connected with strong duality in linear
programming.

Player \(1\) seeks the largest guaranteed payoff.

Player \(2\) seeks the smallest guaranteed upper bound.

Strong duality implies equality of these optimal values.

%%%%%%%%%%%%%%%%%%%%%%%%%%%%%%%%%%%%%%%%%%%%%%%%%%%%%%%%%%%%
\subsection{Constant-Sum Games}
\label{subsec:constant-sum-games}
%%%%%%%%%%%%%%%%%%%%%%%%%%%%%%%%%%%%%%%%%%%%%%%%%%%%%%%%%%%%

A constant-sum game satisfies

\[
\boxed{
u_1(s)+u_2(s)=C
}
\]

for every outcome.

By subtracting constants appropriately, a two-player constant-sum game can
be transformed into a zero-sum game without changing strategic incentives.

%%%%%%%%%%%%%%%%%%%%%%%%%%%%%%%%%%%%%%%%%%%%%%%%%%%%%%%%%%%%
\subsection{General-Sum Games}
\label{subsec:general-sum-games}
%%%%%%%%%%%%%%%%%%%%%%%%%%%%%%%%%%%%%%%%%%%%%%%%%%%%%%%%%%%%

In a general-sum game,

\[
u_1+u_2
\]

need not be constant.

Players may have:

\[
\boxed{
\text{conflicting interests},
}
\]

\[
\boxed{
\text{aligned interests},
}
\]

or both.

Most economic and organizational interactions are general-sum games.

%%%%%%%%%%%%%%%%%%%%%%%%%%%%%%%%%%%%%%%%%%%%%%%%%%%%%%%%%%%%
\subsection{Pareto Efficiency in Games}
\label{subsec:pareto-efficiency-games}
%%%%%%%%%%%%%%%%%%%%%%%%%%%%%%%%%%%%%%%%%%%%%%%%%%%%%%%%%%%%

An outcome is Pareto efficient if no other feasible outcome makes at least
one player better off without making another worse off.

A Nash equilibrium may fail Pareto efficiency.

The Prisoner's Dilemma is the standard example.

%%%%%%%%%%%%%%%%%%%%%%%%%%%%%%%%%%%%%%%%%%%%%%%%%%%%%%%%%%%%
\subsection{Social Welfare}
\label{subsec:social-welfare-games}
%%%%%%%%%%%%%%%%%%%%%%%%%%%%%%%%%%%%%%%%%%%%%%%%%%%%%%%%%%%%

A simple social-welfare measure is

\[
\boxed{
W(s)
=
\sum_i
u_i(s).
}
\]

A social planner may seek

\[
\boxed{
\max_s
W(s).
}
\]

Individual strategic incentives need not produce the welfare-maximizing
outcome.

%%%%%%%%%%%%%%%%%%%%%%%%%%%%%%%%%%%%%%%%%%%%%%%%%%%%%%%%%%%%
\subsection{Price of Anarchy}
\label{subsec:price-of-anarchy}
%%%%%%%%%%%%%%%%%%%%%%%%%%%%%%%%%%%%%%%%%%%%%%%%%%%%%%%%%%%%

The price of anarchy measures efficiency loss caused by decentralized
self-interested behavior.

For a cost-minimization setting, one common definition is

\[
\boxed{
PoA
=
\frac{
\text{worst equilibrium social cost}
}{
\text{minimum possible social cost}
}.
}
\]

Thus

\[
\boxed{
PoA\geq1.
}
\]

Different conventions are used for welfare-maximization problems.

%%%%%%%%%%%%%%%%%%%%%%%%%%%%%%%%%%%%%%%%%%%%%%%%%%%%%%%%%%%%
\subsection{Extensive-Form Games}
\label{subsec:extensive-form-games}
%%%%%%%%%%%%%%%%%%%%%%%%%%%%%%%%%%%%%%%%%%%%%%%%%%%%%%%%%%%%

Extensive-form games explicitly represent:

\[
\boxed{
\text{order of moves},
}
\]

\[
\boxed{
\text{available actions},
}
\]

\[
\boxed{
\text{information},
}
\]

and

\[
\boxed{
\text{terminal payoffs}.
}
\]

They are commonly drawn as game trees.

%%%%%%%%%%%%%%%%%%%%%%%%%%%%%%%%%%%%%%%%%%%%%%%%%%%%%%%%%%%%
\subsection{Decision Nodes in Extensive Form}
\label{subsec:extensive-game-decision-nodes}
%%%%%%%%%%%%%%%%%%%%%%%%%%%%%%%%%%%%%%%%%%%%%%%%%%%%%%%%%%%%

Each decision node belongs to a player.

The branches leaving the node represent available actions.

A terminal node specifies the resulting payoff vector.

%%%%%%%%%%%%%%%%%%%%%%%%%%%%%%%%%%%%%%%%%%%%%%%%%%%%%%%%%%%%
\subsection{Information Sets}
\label{subsec:information-sets}
%%%%%%%%%%%%%%%%%%%%%%%%%%%%%%%%%%%%%%%%%%%%%%%%%%%%%%%%%%%%

An information set contains decision nodes that a player cannot distinguish
when choosing an action.

If a player observes every previous move relevant to their decision, the
game has perfect information.

Simultaneous games can be represented in extensive form using nontrivial
information sets.

%%%%%%%%%%%%%%%%%%%%%%%%%%%%%%%%%%%%%%%%%%%%%%%%%%%%%%%%%%%%
\subsection{Strategy in an Extensive-Form Game}
\label{subsec:extensive-form-strategy}
%%%%%%%%%%%%%%%%%%%%%%%%%%%%%%%%%%%%%%%%%%%%%%%%%%%%%%%%%%%%

A strategy must specify what a player would do at every information set
where that player might be called to act.

Thus an extensive-form strategy is a complete contingent plan.

It includes actions at nodes that may never be reached in equilibrium.

%%%%%%%%%%%%%%%%%%%%%%%%%%%%%%%%%%%%%%%%%%%%%%%%%%%%%%%%%%%%
\subsection{Backward Induction}
\label{subsec:game-backward-induction}
%%%%%%%%%%%%%%%%%%%%%%%%%%%%%%%%%%%%%%%%%%%%%%%%%%%%%%%%%%%%

For a finite game of perfect information, backward induction begins at the
last decision nodes.

At each node:

\[
\boxed{
\text{choose the action maximizing the current player's payoff}.
}
\]

Then move backward through the tree.

This produces strategies that are sequentially rational at every stage.

%%%%%%%%%%%%%%%%%%%%%%%%%%%%%%%%%%%%%%%%%%%%%%%%%%%%%%%%%%%%
\subsection{Worked Example: Backward Induction}
\label{subsec:worked-backward-induction-game}
%%%%%%%%%%%%%%%%%%%%%%%%%%%%%%%%%%%%%%%%%%%%%%%%%%%%%%%%%%%%

\begin{workedexamplebox}{Sequential Entry Decision}

Suppose Player \(1\) chooses:

\[
\text{Enter}
\]

or

\[
\text{Stay Out}.
\]

If Player \(1\) stays out, payoffs are

\[
(0,4).
\]

If Player \(1\) enters, Player \(2\) chooses:

\[
\text{Fight}
\]

or

\[
\text{Accommodate}.
\]

Fight gives

\[
(-2,-1),
\]

while Accommodate gives

\[
(2,2).
\]

At Player \(2\)'s decision node,

\[
2>-1,
\]

so Player \(2\) chooses Accommodate.

Anticipating this, Player \(1\) compares

\[
2
\]

from entering with

\[
0
\]

from staying out.

Therefore,

\begin{answerbox}
\[
\boxed{
\text{Player 1 enters and Player 2 accommodates}.
}
\]
\end{answerbox}

\end{workedexamplebox}

%%%%%%%%%%%%%%%%%%%%%%%%%%%%%%%%%%%%%%%%%%%%%%%%%%%%%%%%%%%%
\subsection{Subgames}
\label{subsec:subgames}
%%%%%%%%%%%%%%%%%%%%%%%%%%%%%%%%%%%%%%%%%%%%%%%%%%%%%%%%%%%%

A subgame is a portion of an extensive-form game beginning at a suitable
decision node and containing all successor nodes without cutting across an
information set.

A subgame can itself be analyzed as a complete strategic game.

%%%%%%%%%%%%%%%%%%%%%%%%%%%%%%%%%%%%%%%%%%%%%%%%%%%%%%%%%%%%
\subsection{Subgame-Perfect Equilibrium}
\label{subsec:subgame-perfect-equilibrium}
%%%%%%%%%%%%%%%%%%%%%%%%%%%%%%%%%%%%%%%%%%%%%%%%%%%%%%%%%%%%

A strategy profile is a subgame-perfect Nash equilibrium if it induces a
Nash equilibrium in every subgame.

Thus it rules out strategies based on threats or promises that would not be
optimal if the relevant decision node were actually reached.

%%%%%%%%%%%%%%%%%%%%%%%%%%%%%%%%%%%%%%%%%%%%%%%%%%%%%%%%%%%%
\subsection{Credible Threats}
\label{subsec:credible-threats}
%%%%%%%%%%%%%%%%%%%%%%%%%%%%%%%%%%%%%%%%%%%%%%%%%%%%%%%%%%%%

A threat is credible if carrying it out would be optimal when the moment
to act arrives.

A threat is noncredible if the threatening player would prefer not to
carry it out once the decision node is reached.

Backward induction eliminates noncredible threats in finite perfect-
information games.

%%%%%%%%%%%%%%%%%%%%%%%%%%%%%%%%%%%%%%%%%%%%%%%%%%%%%%%%%%%%
\subsection{Commitment}
\label{subsec:game-theory-commitment}
%%%%%%%%%%%%%%%%%%%%%%%%%%%%%%%%%%%%%%%%%%%%%%%%%%%%%%%%%%%%

A player may improve their strategic position by committing in advance to
an action.

Commitment changes the other players' expectations and best responses.

Examples include:

\[
\boxed{
\text{capacity investment},
}
\]

\[
\boxed{
\text{contractual commitments},
}
\]

or

\[
\boxed{
\text{automated policies}.
}
\]

%%%%%%%%%%%%%%%%%%%%%%%%%%%%%%%%%%%%%%%%%%%%%%%%%%%%%%%%%%%%
\subsection{First-Mover Advantage}
\label{subsec:first-mover-advantage}
%%%%%%%%%%%%%%%%%%%%%%%%%%%%%%%%%%%%%%%%%%%%%%%%%%%%%%%%%%%%

Moving first can provide an advantage if the first player can commit to a
choice that changes the follower's optimal response.

However, moving first is not always beneficial.

The payoff structure determines whether a first-mover or second-mover
advantage exists.

%%%%%%%%%%%%%%%%%%%%%%%%%%%%%%%%%%%%%%%%%%%%%%%%%%%%%%%%%%%%
\subsection{Stackelberg Competition Preview}
\label{subsec:stackelberg-preview}
%%%%%%%%%%%%%%%%%%%%%%%%%%%%%%%%%%%%%%%%%%%%%%%%%%%%%%%%%%%%

In Stackelberg competition, a leader chooses an action first.

A follower observes that action and responds optimally.

The leader anticipates the follower's reaction.

Thus the leader solves a nested problem of the form

\[
\boxed{
\max_{x}
u_1
(
x,
y^*(x)
),
}
\]

where

\[
\boxed{
y^*(x)
\in
\operatorname*{arg\,max}_y
u_2(x,y).
}
\]

%%%%%%%%%%%%%%%%%%%%%%%%%%%%%%%%%%%%%%%%%%%%%%%%%%%%%%%%%%%%
\subsection{Repeated Games}
\label{subsec:repeated-games}
%%%%%%%%%%%%%%%%%%%%%%%%%%%%%%%%%%%%%%%%%%%%%%%%%%%%%%%%%%%%

A stage game may be played repeatedly.

Repeated interaction can change incentives because current actions can
affect future behavior.

A repeated-game strategy may therefore depend on the entire observed
history.

%%%%%%%%%%%%%%%%%%%%%%%%%%%%%%%%%%%%%%%%%%%%%%%%%%%%%%%%%%%%
\subsection{Finite Repetition}
\label{subsec:finite-repeated-games}
%%%%%%%%%%%%%%%%%%%%%%%%%%%%%%%%%%%%%%%%%%%%%%%%%%%%%%%%%%%%

Suppose a stage game has a unique Nash equilibrium and is repeated a known
finite number of times.

Backward induction may cause the stage-game equilibrium to be played in
every period under standard assumptions.

The last period has no future punishment or reward.

This logic then propagates backward.

%%%%%%%%%%%%%%%%%%%%%%%%%%%%%%%%%%%%%%%%%%%%%%%%%%%%%%%%%%%%
\subsection{Infinite Repetition}
\label{subsec:infinite-repeated-games}
%%%%%%%%%%%%%%%%%%%%%%%%%%%%%%%%%%%%%%%%%%%%%%%%%%%%%%%%%%%%

With indefinite or infinite repetition, future consequences can sustain
cooperation that would not be stable in a one-shot game.

A discounted payoff may be

\[
\boxed{
\sum_{t=0}^{\infty}
\delta^t
u_i^t,
}
\]

where

\[
\boxed{
0<\delta<1.
}
\]

Larger \(\delta\) places more weight on future payoffs.

%%%%%%%%%%%%%%%%%%%%%%%%%%%%%%%%%%%%%%%%%%%%%%%%%%%%%%%%%%%%
\subsection{Trigger Strategies}
\label{subsec:trigger-strategies}
%%%%%%%%%%%%%%%%%%%%%%%%%%%%%%%%%%%%%%%%%%%%%%%%%%%%%%%%%%%%

A trigger strategy rewards cooperation while threatening future punishment
after deviation.

For example, a grim-trigger strategy may:

\begin{enumerate}
    \item cooperate while all players have cooperated;
    \item after any deviation, switch permanently to a punishment strategy.
\end{enumerate}

The threat can deter deviation when future losses are sufficiently large.

%%%%%%%%%%%%%%%%%%%%%%%%%%%%%%%%%%%%%%%%%%%%%%%%%%%%%%%%%%%%
\subsection{Folk Theorem Preview}
\label{subsec:folk-theorem-preview}
%%%%%%%%%%%%%%%%%%%%%%%%%%%%%%%%%%%%%%%%%%%%%%%%%%%%%%%%%%%%

Folk theorems show that repeated games can support many payoff outcomes as
equilibria when players are sufficiently patient and appropriate
punishments are available.

Thus repetition can produce substantially more equilibrium behavior than
the one-shot game.

Exact statements depend on the repeated-game framework.

%%%%%%%%%%%%%%%%%%%%%%%%%%%%%%%%%%%%%%%%%%%%%%%%%%%%%%%%%%%%
\subsection{Games of Incomplete Information}
\label{subsec:incomplete-information-games}
%%%%%%%%%%%%%%%%%%%%%%%%%%%%%%%%%%%%%%%%%%%%%%%%%%%%%%%%%%%%

In a game of incomplete information, some player does not know all relevant
characteristics of another player.

Examples include uncertainty about:

\[
\boxed{
\text{cost},
}
\]

\[
\boxed{
\text{valuation},
}
\]

\[
\boxed{
\text{product quality},
}
\]

or

\[
\boxed{
\text{risk type}.
}
\]

%%%%%%%%%%%%%%%%%%%%%%%%%%%%%%%%%%%%%%%%%%%%%%%%%%%%%%%%%%%%
\subsection{Types}
\label{subsec:game-theory-types}
%%%%%%%%%%%%%%%%%%%%%%%%%%%%%%%%%%%%%%%%%%%%%%%%%%%%%%%%%%%%

A player's private information is represented by a type.

Player \(i\)'s type is

\[
\boxed{
\theta_i\in\Theta_i.
}
\]

Payoffs may depend on the full type profile:

\[
\boxed{
u_i
(
s_1,\ldots,s_n,
\theta_1,\ldots,\theta_n
).
}
\]

%%%%%%%%%%%%%%%%%%%%%%%%%%%%%%%%%%%%%%%%%%%%%%%%%%%%%%%%%%%%
\subsection{Bayesian Games}
\label{subsec:bayesian-games}
%%%%%%%%%%%%%%%%%%%%%%%%%%%%%%%%%%%%%%%%%%%%%%%%%%%%%%%%%%%%

A Bayesian game specifies:

\[
\boxed{
\text{players},
}
\]

\[
\boxed{
\text{types},
}
\]

\[
\boxed{
\text{beliefs},
}
\]

\[
\boxed{
\text{strategies},
}
\]

and

\[
\boxed{
\text{type-dependent payoffs}.
}
\]

A strategy maps a player's type to an action.

%%%%%%%%%%%%%%%%%%%%%%%%%%%%%%%%%%%%%%%%%%%%%%%%%%%%%%%%%%%%
\subsection{Bayesian Nash Equilibrium}
\label{subsec:bayesian-nash-equilibrium}
%%%%%%%%%%%%%%%%%%%%%%%%%%%%%%%%%%%%%%%%%%%%%%%%%%%%%%%%%%%%

In Bayesian Nash equilibrium, each player's strategy maximizes expected
payoff given:

\[
\boxed{
\text{their own type},
}
\]

\[
\boxed{
\text{their beliefs about others' types},
}
\]

and

\[
\boxed{
\text{the equilibrium strategies of the other players}.
}
\]

This extends Nash equilibrium to incomplete information.

%%%%%%%%%%%%%%%%%%%%%%%%%%%%%%%%%%%%%%%%%%%%%%%%%%%%%%%%%%%%
\subsection{Signaling}
\label{subsec:game-theory-signaling}
%%%%%%%%%%%%%%%%%%%%%%%%%%%%%%%%%%%%%%%%%%%%%%%%%%%%%%%%%%%%

A signaling game occurs when an informed player takes an observable action
that may reveal information about their private type.

Examples include:

\[
\boxed{
\text{education signaling productivity},
}
\]

\[
\boxed{
\text{warranties signaling product quality},
}
\]

or

\[
\boxed{
\text{prices signaling market information}.
}
\]

%%%%%%%%%%%%%%%%%%%%%%%%%%%%%%%%%%%%%%%%%%%%%%%%%%%%%%%%%%%%
\subsection{Separating and Pooling Equilibria}
\label{subsec:separating-pooling}
%%%%%%%%%%%%%%%%%%%%%%%%%%%%%%%%%%%%%%%%%%%%%%%%%%%%%%%%%%%%

In a separating equilibrium, different types choose different observable
actions.

Observers can infer type from the signal.

In a pooling equilibrium, multiple types choose the same signal.

The observer cannot fully infer type from the observed action.

%%%%%%%%%%%%%%%%%%%%%%%%%%%%%%%%%%%%%%%%%%%%%%%%%%%%%%%%%%%%
\subsection{Screening}
\label{subsec:game-theory-screening}
%%%%%%%%%%%%%%%%%%%%%%%%%%%%%%%%%%%%%%%%%%%%%%%%%%%%%%%%%%%%

Screening occurs when an uninformed party designs choices that induce
different types to reveal themselves through their selections.

Examples include:

\[
\boxed{
\text{insurance contracts},
}
\]

\[
\boxed{
\text{pricing plans},
}
\]

and

\[
\boxed{
\text{employment contracts}.
}
\]

%%%%%%%%%%%%%%%%%%%%%%%%%%%%%%%%%%%%%%%%%%%%%%%%%%%%%%%%%%%%
\subsection{Auctions}
\label{subsec:game-theory-auctions}
%%%%%%%%%%%%%%%%%%%%%%%%%%%%%%%%%%%%%%%%%%%%%%%%%%%%%%%%%%%%

An auction is a game in which bidders compete for one or more items.

Important formats include:

\[
\boxed{
\text{English auction},
}
\]

\[
\boxed{
\text{Dutch auction},
}
\]

\[
\boxed{
\text{first-price sealed-bid auction},
}
\]

and

\[
\boxed{
\text{second-price sealed-bid auction}.
}
\]

%%%%%%%%%%%%%%%%%%%%%%%%%%%%%%%%%%%%%%%%%%%%%%%%%%%%%%%%%%%%
\subsection{Private Values}
\label{subsec:private-value-auctions}
%%%%%%%%%%%%%%%%%%%%%%%%%%%%%%%%%%%%%%%%%%%%%%%%%%%%%%%%%%%%

In a private-value auction, bidder \(i\) knows their own valuation

\[
\boxed{
v_i
}
\]

for the object.

Other bidders' valuations are uncertain.

A bidder's utility may be

\[
\boxed{
v_i-\text{payment}
}
\]

if they win, and

\[
\boxed{
0
}
\]

if they lose.

%%%%%%%%%%%%%%%%%%%%%%%%%%%%%%%%%%%%%%%%%%%%%%%%%%%%%%%%%%%%
\subsection{First-Price Sealed-Bid Auction}
\label{subsec:first-price-auction}
%%%%%%%%%%%%%%%%%%%%%%%%%%%%%%%%%%%%%%%%%%%%%%%%%%%%%%%%%%%%

Each bidder submits one sealed bid.

The highest bidder wins and pays their own bid.

Thus winning payoff is

\[
\boxed{
v_i-b_i.
}
\]

Bidders generally have an incentive to shade their bids below their true
valuation under standard private-value assumptions.

%%%%%%%%%%%%%%%%%%%%%%%%%%%%%%%%%%%%%%%%%%%%%%%%%%%%%%%%%%%%
\subsection{Second-Price Sealed-Bid Auction}
\label{subsec:second-price-auction}
%%%%%%%%%%%%%%%%%%%%%%%%%%%%%%%%%%%%%%%%%%%%%%%%%%%%%%%%%%%%

In a second-price sealed-bid auction:

\[
\boxed{
\text{highest bidder wins},
}
\]

but pays

\[
\boxed{
\text{the second-highest bid}.
}
\]

This is also called a Vickrey auction.

%%%%%%%%%%%%%%%%%%%%%%%%%%%%%%%%%%%%%%%%%%%%%%%%%%%%%%%%%%%%
\subsection{Truthful Bidding in a Second-Price Auction}
\label{subsec:truthful-vickrey}
%%%%%%%%%%%%%%%%%%%%%%%%%%%%%%%%%%%%%%%%%%%%%%%%%%%%%%%%%%%%

Under the standard independent private-value model, bidding one's true
valuation is a weakly dominant strategy.

Suppose bidder \(i\) has value \(v_i\).

Their bid affects whether they win, but conditional on winning, the price
is determined by the highest competing bid.

Thus exaggerating or understating \(v_i\) cannot improve payoff across all
possible competing bids.

%%%%%%%%%%%%%%%%%%%%%%%%%%%%%%%%%%%%%%%%%%%%%%%%%%%%%%%%%%%%
\subsection{Mechanism Design Preview}
\label{subsec:mechanism-design-preview}
%%%%%%%%%%%%%%%%%%%%%%%%%%%%%%%%%%%%%%%%%%%%%%%%%%%%%%%%%%%%

Game theory analyzes behavior under given rules.

Mechanism design asks:

\[
\boxed{
\text{What rules should be designed to produce desirable strategic
behavior?}
}
\]

The designer chooses an institution, auction, contract, or allocation rule
while anticipating equilibrium responses.

%%%%%%%%%%%%%%%%%%%%%%%%%%%%%%%%%%%%%%%%%%%%%%%%%%%%%%%%%%%%
\subsection{Incentive Compatibility}
\label{subsec:incentive-compatibility}
%%%%%%%%%%%%%%%%%%%%%%%%%%%%%%%%%%%%%%%%%%%%%%%%%%%%%%%%%%%%

A mechanism is incentive compatible when agents find it optimal to behave
in the intended way.

For example, truthful reporting may satisfy

\[
\boxed{
\text{truth telling}
\in
\operatorname*{arg\,max}
\text{ expected utility}.
}
\]

Incentive constraints are central to mechanism design.

%%%%%%%%%%%%%%%%%%%%%%%%%%%%%%%%%%%%%%%%%%%%%%%%%%%%%%%%%%%%
\subsection{Congestion Games}
\label{subsec:congestion-games}
%%%%%%%%%%%%%%%%%%%%%%%%%%%%%%%%%%%%%%%%%%%%%%%%%%%%%%%%%%%%

In a congestion game, players choose shared resources.

A resource's cost increases as more players use it.

Examples include:

\[
\boxed{
\text{road routes},
}
\]

\[
\boxed{
\text{communication channels},
}
\]

and

\[
\boxed{
\text{shared computing resources}.
}
\]

%%%%%%%%%%%%%%%%%%%%%%%%%%%%%%%%%%%%%%%%%%%%%%%%%%%%%%%%%%%%
\subsection{Routing Game Example}
\label{subsec:routing-game-example}
%%%%%%%%%%%%%%%%%%%%%%%%%%%%%%%%%%%%%%%%%%%%%%%%%%%%%%%%%%%%

Suppose travelers choose between two routes.

Route costs depend on flow:

\[
\boxed{
c_1(x_1),
\qquad
c_2(x_2).
}
\]

Each traveler prefers a lower-cost route.

Yet individually optimal route choices can create collectively inefficient
congestion.

%%%%%%%%%%%%%%%%%%%%%%%%%%%%%%%%%%%%%%%%%%%%%%%%%%%%%%%%%%%%
\subsection{Wardrop Equilibrium Preview}
\label{subsec:wardrop-equilibrium-preview}
%%%%%%%%%%%%%%%%%%%%%%%%%%%%%%%%%%%%%%%%%%%%%%%%%%%%%%%%%%%%

In a nonatomic routing model, a Wardrop equilibrium has the property that
all used routes between an origin and destination have equal minimal
travel cost.

No infinitesimal traveler can improve travel time by switching routes
unilaterally.

This is a continuous-population analogue of Nash equilibrium.

%%%%%%%%%%%%%%%%%%%%%%%%%%%%%%%%%%%%%%%%%%%%%%%%%%%%%%%%%%%%
\subsection{Potential Games}
\label{subsec:potential-games}
%%%%%%%%%%%%%%%%%%%%%%%%%%%%%%%%%%%%%%%%%%%%%%%%%%%%%%%%%%%%

A potential game has a scalar function

\[
\boxed{
\Phi(s)
}
\]

that tracks incentives for unilateral deviations.

For an exact potential game,

\[
\boxed{
u_i(s_i',s_{-i})
-
u_i(s_i,s_{-i})
=
\Phi(s_i',s_{-i})
-
\Phi(s_i,s_{-i}).
}
\]

Thus unilateral payoff improvements correspond to increases in the
potential.

%%%%%%%%%%%%%%%%%%%%%%%%%%%%%%%%%%%%%%%%%%%%%%%%%%%%%%%%%%%%
\subsection{Existence in Finite Potential Games}
\label{subsec:potential-game-existence}
%%%%%%%%%%%%%%%%%%%%%%%%%%%%%%%%%%%%%%%%%%%%%%%%%%%%%%%%%%%%

A finite exact potential game has at least one pure-strategy Nash
equilibrium.

Because there are finitely many strategy profiles, the potential function
has a maximizer.

No unilateral deviation can increase the player's payoff at such a
maximizer.

%%%%%%%%%%%%%%%%%%%%%%%%%%%%%%%%%%%%%%%%%%%%%%%%%%%%%%%%%%%%
\subsection{Evolutionary Game Theory}
\label{subsec:evolutionary-game-theory}
%%%%%%%%%%%%%%%%%%%%%%%%%%%%%%%%%%%%%%%%%%%%%%%%%%%%%%%%%%%%

Evolutionary game theory studies populations in which strategies reproduce,
spread, or disappear according to relative success.

Players need not be perfectly rational.

Instead, successful strategies may become more common over time.

%%%%%%%%%%%%%%%%%%%%%%%%%%%%%%%%%%%%%%%%%%%%%%%%%%%%%%%%%%%%
\subsection{Population Shares}
\label{subsec:evolutionary-population-shares}
%%%%%%%%%%%%%%%%%%%%%%%%%%%%%%%%%%%%%%%%%%%%%%%%%%%%%%%%%%%%

Suppose strategy \(i\) has population share

\[
\boxed{
x_i\geq0,
}
\]

with

\[
\boxed{
\sum_i
x_i=1.
}
\]

The expected payoff to strategy \(i\) may depend on the population state
\(x\).

%%%%%%%%%%%%%%%%%%%%%%%%%%%%%%%%%%%%%%%%%%%%%%%%%%%%%%%%%%%%
\subsection{Replicator Dynamics}
\label{subsec:replicator-dynamics}
%%%%%%%%%%%%%%%%%%%%%%%%%%%%%%%%%%%%%%%%%%%%%%%%%%%%%%%%%%%%

A standard evolutionary model is

\[
\boxed{
\dot{x}_i
=
x_i
\left(
u_i(x)
-
\bar u(x)
\right),
}
\]

where

\[
\boxed{
\bar u(x)
=
\sum_j
x_j u_j(x)
}
\]

is average population payoff.

Strategies performing above average increase in frequency.

%%%%%%%%%%%%%%%%%%%%%%%%%%%%%%%%%%%%%%%%%%%%%%%%%%%%%%%%%%%%
\subsection{Evolutionarily Stable Strategy}
\label{subsec:evolutionarily-stable-strategy}
%%%%%%%%%%%%%%%%%%%%%%%%%%%%%%%%%%%%%%%%%%%%%%%%%%%%%%%%%%%%

An evolutionarily stable strategy, abbreviated ESS, is a strategy that
cannot be successfully invaded by a sufficiently small population using an
alternative strategy.

ESS is stronger than merely being a symmetric Nash equilibrium in many
settings.

%%%%%%%%%%%%%%%%%%%%%%%%%%%%%%%%%%%%%%%%%%%%%%%%%%%%%%%%%%%%
\subsection{Game Theory and Optimization}
\label{subsec:game-theory-optimization-connection}
%%%%%%%%%%%%%%%%%%%%%%%%%%%%%%%%%%%%%%%%%%%%%%%%%%%%%%%%%%%%

Ordinary optimization asks one decision maker to solve

\[
\boxed{
\max_x
f(x).
}
\]

Game theory asks several players to solve coupled problems:

\[
\boxed{
\max_{x_i}
u_i(x_i,x_{-i}),
\qquad
i=1,\ldots,n.
}
\]

Equilibrium requires these optimizations to be mutually consistent.

%%%%%%%%%%%%%%%%%%%%%%%%%%%%%%%%%%%%%%%%%%%%%%%%%%%%%%%%%%%%
\subsection{Game Theory and Variational Inequalities Preview}
\label{subsec:game-theory-variational-inequalities}
%%%%%%%%%%%%%%%%%%%%%%%%%%%%%%%%%%%%%%%%%%%%%%%%%%%%%%%%%%%%

Some continuous games can be represented as variational inequalities.

If players' first-order conditions are combined into mapping \(F(x)\), an
equilibrium may satisfy

\[
\boxed{
F(x^*)^T
(x-x^*)
\geq
0
}
\]

for all feasible \(x\).

This connects equilibrium computation with mathematical optimization.

%%%%%%%%%%%%%%%%%%%%%%%%%%%%%%%%%%%%%%%%%%%%%%%%%%%%%%%%%%%%
\subsection{Game Theory and Decision Theory}
\label{subsec:game-theory-decision-connection}
%%%%%%%%%%%%%%%%%%%%%%%%%%%%%%%%%%%%%%%%%%%%%%%%%%%%%%%%%%%%

Decision theory studies choice under uncertainty.

Game theory studies choice under strategic interaction.

A key difference is:

\[
\boxed{
\text{nature does not optimize against you},
}
\]

while

\[
\boxed{
\text{another player may}.
}
\]

Thus beliefs about other players must account for their incentives.

%%%%%%%%%%%%%%%%%%%%%%%%%%%%%%%%%%%%%%%%%%%%%%%%%%%%%%%%%%%%
\subsection{Game Theory and Multiobjective Optimization}
\label{subsec:game-theory-multiobjective}
%%%%%%%%%%%%%%%%%%%%%%%%%%%%%%%%%%%%%%%%%%%%%%%%%%%%%%%%%%%%

Multiobjective optimization represents several objectives for a decision
maker or planning system.

Game theory represents several decision makers with potentially different
objectives.

A social planner might optimize a vector of player utilities, while
strategic equilibrium asks what happens when players optimize
independently.

%%%%%%%%%%%%%%%%%%%%%%%%%%%%%%%%%%%%%%%%%%%%%%%%%%%%%%%%%%%%
\subsection{Game Theory and Networks}
\label{subsec:game-theory-networks}
%%%%%%%%%%%%%%%%%%%%%%%%%%%%%%%%%%%%%%%%%%%%%%%%%%%%%%%%%%%%

Network decisions can be strategic.

Examples include:

\[
\boxed{
\text{selfish routing},
}
\]

\[
\boxed{
\text{network formation},
}
\]

\[
\boxed{
\text{bandwidth sharing},
}
\]

and

\[
\boxed{
\text{security investment}.
}
\]

Network optimization asks what is centrally optimal.

Network games ask what decentralized players choose.

%%%%%%%%%%%%%%%%%%%%%%%%%%%%%%%%%%%%%%%%%%%%%%%%%%%%%%%%%%%%
\subsection{Game Theory Workflow}
\label{subsec:game-theory-workflow}
%%%%%%%%%%%%%%%%%%%%%%%%%%%%%%%%%%%%%%%%%%%%%%%%%%%%%%%%%%%%

A practical workflow is:

\begin{enumerate}
    \item identify the players;
    \item identify each player's feasible strategies;
    \item define payoffs;
    \item determine whether moves are simultaneous or sequential;
    \item determine what information each player observes;
    \item identify dominant and dominated strategies;
    \item compute best responses;
    \item search for pure Nash equilibria;
    \item if necessary, compute mixed equilibria;
    \item use backward induction for finite sequential games;
    \item check subgame perfection;
    \item incorporate beliefs when information is incomplete;
    \item compare equilibrium outcomes with Pareto-efficient or socially
          optimal outcomes;
    \item test sensitivity to payoff assumptions;
    \item interpret whether commitment, information, or institutional
          changes alter equilibrium incentives.
\end{enumerate}

%%%%%%%%%%%%%%%%%%%%%%%%%%%%%%%%%%%%%%%%%%%%%%%%%%%%%%%%%%%%
\subsection{Choosing a Game-Theoretic Method}
\label{subsec:choosing-game-method}
%%%%%%%%%%%%%%%%%%%%%%%%%%%%%%%%%%%%%%%%%%%%%%%%%%%%%%%%%%%%

For a finite simultaneous game, use:

\[
\boxed{
\text{best responses and Nash equilibrium}.
}
\]

For a two-player zero-sum game, use:

\[
\boxed{
\text{minimax analysis or linear programming}.
}
\]

For a finite sequential perfect-information game, use:

\[
\boxed{
\text{backward induction}.
}
\]

For sequential games with strategic credibility, use:

\[
\boxed{
\text{subgame-perfect equilibrium}.
}
\]

For incomplete information, use:

\[
\boxed{
\text{Bayesian games}.
}
\]

For repeated interaction, use:

\[
\boxed{
\text{repeated-game equilibrium analysis}.
}
\]

%%%%%%%%%%%%%%%%%%%%%%%%%%%%%%%%%%%%%%%%%%%%%%%%%%%%%%%%%%%%
\subsection{Common Errors}
\label{subsec:game-theory-common-errors}
%%%%%%%%%%%%%%%%%%%%%%%%%%%%%%%%%%%%%%%%%%%%%%%%%%%%%%%%%%%%

\subsubsection{Confusing Strategy and Action}

An action is a choice at a particular decision point.

A strategy is a complete plan specifying actions at every relevant
contingency.

In simultaneous one-shot games, the distinction may appear invisible.

%%%%%%%%%%%%%%%%%%%%%%%%%%%%%%%%%%%%%%%%%%%%%%%%%%%%%%%%%%%%
\subsubsection{Confusing Dominant Strategy and Best Response}

A dominant strategy is optimal against every opponent strategy.

A best response is optimal only against a specified opponent strategy.

%%%%%%%%%%%%%%%%%%%%%%%%%%%%%%%%%%%%%%%%%%%%%%%%%%%%%%%%%%%%
\subsubsection{Assuming Every Game Has a Pure Nash Equilibrium}

Finite games always have a mixed-strategy Nash equilibrium.

They do not necessarily have a pure-strategy equilibrium.

Matching Pennies is a standard example.

%%%%%%%%%%%%%%%%%%%%%%%%%%%%%%%%%%%%%%%%%%%%%%%%%%%%%%%%%%%%
\subsubsection{Assuming Nash Equilibrium Maximizes Total Welfare}

A Nash equilibrium only prevents profitable unilateral deviations.

It does not necessarily maximize

\[
\sum_i
u_i.
\]

%%%%%%%%%%%%%%%%%%%%%%%%%%%%%%%%%%%%%%%%%%%%%%%%%%%%%%%%%%%%
\subsubsection{Confusing Pareto Efficiency and Nash Equilibrium}

Pareto efficiency is a welfare concept.

Nash equilibrium is an incentive-stability concept.

An outcome can satisfy one without satisfying the other.

%%%%%%%%%%%%%%%%%%%%%%%%%%%%%%%%%%%%%%%%%%%%%%%%%%%%%%%%%%%%
\subsubsection{Ignoring Mixed Strategies}

When no pure equilibrium exists, randomization may be essential.

Checking only payoff-matrix cells can therefore miss equilibrium entirely.

%%%%%%%%%%%%%%%%%%%%%%%%%%%%%%%%%%%%%%%%%%%%%%%%%%%%%%%%%%%%
\subsubsection{Making a Player Indifferent to Their Own Mixture}

To determine Player \(1\)'s equilibrium mixing probability, usually make
Player \(2\) indifferent among the pure strategies used in support.

Likewise, determine Player \(2\)'s mixture by making Player \(1\)
indifferent.

%%%%%%%%%%%%%%%%%%%%%%%%%%%%%%%%%%%%%%%%%%%%%%%%%%%%%%%%%%%%
\subsubsection{Using Maximin as a General Nash-Equilibrium Method}

Maximin is especially natural in zero-sum games.

In general-sum games, Nash equilibrium is based on mutual best responses,
not worst-case payoff alone.

%%%%%%%%%%%%%%%%%%%%%%%%%%%%%%%%%%%%%%%%%%%%%%%%%%%%%%%%%%%%
\subsubsection{Ignoring Off-Path Actions in Sequential Games}

A strategy must specify behavior even at decision nodes that are not
reached along the equilibrium path.

These off-path actions matter for credibility and subgame perfection.

%%%%%%%%%%%%%%%%%%%%%%%%%%%%%%%%%%%%%%%%%%%%%%%%%%%%%%%%%%%%
\subsubsection{Accepting Noncredible Threats}

A threat that would not be optimal to carry out once reached should not
support a subgame-perfect equilibrium.

Backward induction reveals such threats.

%%%%%%%%%%%%%%%%%%%%%%%%%%%%%%%%%%%%%%%%%%%%%%%%%%%%%%%%%%%%
\subsubsection{Treating Incomplete Information as Random Nature Only}

In Bayesian games, unknown information belongs to strategic players whose
actions depend on their private types.

This differs from a passive state-of-nature model.

%%%%%%%%%%%%%%%%%%%%%%%%%%%%%%%%%%%%%%%%%%%%%%%%%%%%%%%%%%%%
\subsubsection{Assuming Bidding Your Value Is Optimal in Every Auction}

Truthful bidding is weakly dominant in the standard second-price
private-value auction.

It is not generally the equilibrium strategy in a first-price auction.

%%%%%%%%%%%%%%%%%%%%%%%%%%%%%%%%%%%%%%%%%%%%%%%%%%%%%%%%%%%%
\subsubsection{Ignoring Equilibrium Multiplicity}

A game may have several Nash equilibria.

Additional considerations may be needed to predict which equilibrium will
occur.

%%%%%%%%%%%%%%%%%%%%%%%%%%%%%%%%%%%%%%%%%%%%%%%%%%%%%%%%%%%%
\subsubsection{Ignoring Information Structure}

Who knows what, and when, can completely change the equilibrium.

The same payoff possibilities can produce different strategic behavior
under different information structures.

%%%%%%%%%%%%%%%%%%%%%%%%%%%%%%%%%%%%%%%%%%%%%%%%%%%%%%%%%%%%
\subsubsection{Assuming Repetition Always Creates Cooperation}

Repeated interaction can support cooperation under suitable incentives,
but cooperation is not automatic.

Discount factors, monitoring, punishment strategies, and equilibrium
requirements matter.

%%%%%%%%%%%%%%%%%%%%%%%%%%%%%%%%%%%%%%%%%%%%%%%%%%%%%%%%%%%%
\subsection{Applications}
\label{subsec:game-theory-applications}
%%%%%%%%%%%%%%%%%%%%%%%%%%%%%%%%%%%%%%%%%%%%%%%%%%%%%%%%%%%%

\begin{applicationbox}

\textbf{Economics.}
Firms choose prices, quantities, capacity, advertising, and market-entry
strategies while anticipating competitors.

\medskip

\textbf{Auctions and Markets.}
Bidders strategically choose bids, while market designers choose rules to
shape incentives.

\medskip

\textbf{Transportation.}
Drivers choosing routes independently can create congestion and socially
inefficient traffic patterns.

\medskip

\textbf{Networks and Computing.}
Distributed agents compete for communication bandwidth, processing
capacity, storage, and shared resources.

\medskip

\textbf{Operations Research.}
Supply-chain partners, competing firms, service providers, and decentralized
decision makers may pursue different objectives.

\medskip

\textbf{Biology.}
Evolutionary game theory models competition among behavioral or biological
strategies.

\medskip

\textbf{Negotiation and Bargaining.}
Game-theoretic models analyze how strategic parties divide gains from
cooperation.

\medskip

\textbf{Public Policy.}
Regulation can change incentives and therefore alter equilibrium behavior
rather than merely impose technical constraints.

\end{applicationbox}

%%%%%%%%%%%%%%%%%%%%%%%%%%%%%%%%%%%%%%%%%%%%%%%%%%%%%%%%%%%%
\subsection{Exercises}
\label{subsec:game-theory-exercises}
%%%%%%%%%%%%%%%%%%%%%%%%%%%%%%%%%%%%%%%%%%%%%%%%%%%%%%%%%%%%

\begin{exercise}[1]
Define a player.
\end{exercise}

\begin{exercise}[1]
Define a strategy.
\end{exercise}

\begin{exercise}[1]
Define a payoff function.
\end{exercise}

\begin{exercise}[1]
Define a best response.
\end{exercise}

\begin{exercise}[1]
Define Nash equilibrium.
\end{exercise}

\begin{exercise}[1]
Define a dominant strategy.
\end{exercise}

\begin{exercise}[1]
Define a mixed strategy.
\end{exercise}

\begin{exercise}[1]
Define a zero-sum game.
\end{exercise}

\begin{exercise}[1]
Define a subgame-perfect equilibrium.
\end{exercise}

\begin{exercise}[1]
Define a Bayesian game conceptually.
\end{exercise}

\begin{exercise}[2]
Identify all strictly dominated strategies in a given payoff matrix.
\end{exercise}

\begin{exercise}[2]
Perform iterated elimination of strictly dominated strategies.
\end{exercise}

\begin{exercise}[2]
Mark each player's best responses in a \(2\times2\) game.
\end{exercise}

\begin{exercise}[2]
Identify all pure-strategy Nash equilibria.
\end{exercise}

\begin{exercise}[2]
Explain why a dominant-strategy equilibrium is also a Nash equilibrium.
\end{exercise}

\begin{exercise}[2]
Explain why the converse need not hold.
\end{exercise}

\begin{exercise}[2]
Analyze a Prisoner's Dilemma payoff matrix.
\end{exercise}

\begin{exercise}[2]
Identify the pure equilibria of a coordination game.
\end{exercise}

\begin{exercise}[2]
Identify the pure equilibria of Battle of the Sexes.
\end{exercise}

\begin{exercise}[2]
Explain why Matching Pennies has no pure Nash equilibrium.
\end{exercise}

\begin{exercise}[3]
Compute the mixed-strategy equilibrium of Matching Pennies.
\end{exercise}

\begin{exercise}[3]
For a general \(2\times2\) game, derive the probability that makes
Player \(1\) indifferent between two pure strategies.
\end{exercise}

\begin{exercise}[3]
Derive the corresponding equilibrium probability for Player \(1\).
\end{exercise}

\begin{exercise}[3]
Verify that all pure strategies in equilibrium support give equal expected
payoff.
\end{exercise}

\begin{exercise}[3]
Calculate the expected payoff of two mixed strategies using

\[
p^TAq.
\]
\end{exercise}

\begin{exercise}[3]
Compute the pure maximin and minimax values of a zero-sum matrix game.
\end{exercise}

\begin{exercise}[3]
Determine whether the game has a saddle point.
\end{exercise}

\begin{exercise}[3]
Formulate Player \(1\)'s zero-sum optimization problem as a linear program.
\end{exercise}

\begin{exercise}[3]
Formulate Player \(2\)'s corresponding linear program.
\end{exercise}

\begin{exercise}[3]
Interpret the value of a zero-sum game.
\end{exercise}

\begin{exercise}[3]
Construct an extensive-form game from a short sequential scenario.
\end{exercise}

\begin{exercise}[3]
Solve the previous game using backward induction.
\end{exercise}

\begin{exercise}[3]
Identify a noncredible threat in a sequential game.
\end{exercise}

\begin{exercise}[3]
Determine a subgame-perfect equilibrium.
\end{exercise}

\begin{exercise}[3]
Compare the Nash equilibria of the normal form with the subgame-perfect
equilibria of the extensive form.
\end{exercise}

\begin{exercise}[3]
Construct a simple entry-deterrence game.
\end{exercise}

\begin{exercise}[3]
Explain how commitment can change the entry-deterrence outcome.
\end{exercise}

\begin{exercise}[3]
Compute a discounted payoff

\[
\sum_{t=0}^{\infty}
\delta^tu
\]

for a constant stage payoff \(u\).
\end{exercise}

\begin{exercise}[3]
Derive a simple condition under which future punishment deters a one-time
deviation in a repeated Prisoner's Dilemma.
\end{exercise}

\begin{exercise}[3]
Explain the difference between finite and infinite repetition.
\end{exercise}

\begin{exercise}[3]
Construct a Bayesian game with two possible types for one player.
\end{exercise}

\begin{exercise}[3]
Write a type-contingent strategy.
\end{exercise}

\begin{exercise}[3]
Explain the difference between pooling and separating equilibria.
\end{exercise}

\begin{exercise}[3]
Explain the difference between signaling and screening.
\end{exercise}

\begin{exercise}[3]
Compare incentives in first-price and second-price sealed-bid auctions.
\end{exercise}

\begin{exercise}[3]
Explain why truthful bidding is weakly dominant in a standard
second-price auction.
\end{exercise}

\begin{exercise}[3]
Construct a small congestion game.
\end{exercise}

\begin{exercise}[3]
Identify all pure Nash equilibria in the congestion game.
\end{exercise}

\begin{exercise}[3]
Compute the social optimum and compare it with equilibrium.
\end{exercise}

\begin{exercise}[3]
Compute a price of anarchy for a small cost-minimization game.
\end{exercise}

\begin{exercise}[3]
Construct a finite exact potential game.
\end{exercise}

\begin{exercise}[3]
Verify that unilateral payoff differences equal potential differences.
\end{exercise}

\begin{exercise}[3]
Write the replicator equation for a two-strategy population.
\end{exercise}

\begin{exercise}[3]
Find stationary population shares for a simple evolutionary game.
\end{exercise}

\begin{exercise}[4]
Prove that every dominant-strategy equilibrium is a Nash equilibrium.
\end{exercise}

\begin{exercise}[4]
Study rationalizability and iterated elimination of dominated strategies.
\end{exercise}

\begin{exercise}[4]
Study existence of mixed Nash equilibrium in finite games.
\end{exercise}

\begin{exercise}[4]
Prove the minimax theorem using linear programming duality.
\end{exercise}

\begin{exercise}[4]
Study Lemke--Howson methods for bimatrix games.
\end{exercise}

\begin{exercise}[4]
Study correlated equilibrium.
\end{exercise}

\begin{exercise}[4]
Compare correlated equilibrium with Nash equilibrium.
\end{exercise}

\begin{exercise}[4]
Study backward induction and subgame-perfect equilibrium formally.
\end{exercise}

\begin{exercise}[4]
Study extensive games with imperfect information.
\end{exercise}

\begin{exercise}[4]
Study sequential equilibrium conceptually.
\end{exercise}

\begin{exercise}[4]
Study trembling-hand perfection conceptually.
\end{exercise}

\begin{exercise}[4]
Analyze finitely repeated Prisoner's Dilemma.
\end{exercise}

\begin{exercise}[4]
Study Folk Theorem results for infinitely repeated games.
\end{exercise}

\begin{exercise}[4]
Derive Bayesian Nash equilibrium in a simple incomplete-information game.
\end{exercise}

\begin{exercise}[4]
Study signaling equilibria.
\end{exercise}

\begin{exercise}[4]
Study screening and incentive compatibility.
\end{exercise}

\begin{exercise}[4]
Derive symmetric equilibrium bidding in a first-price private-value
auction.
\end{exercise}

\begin{exercise}[4]
Prove truthful bidding in a Vickrey auction.
\end{exercise}

\begin{exercise}[4]
Study the revelation principle conceptually.
\end{exercise}

\begin{exercise}[4]
Study Vickrey--Clarke--Groves mechanisms conceptually.
\end{exercise}

\begin{exercise}[4]
Study Wardrop equilibria in nonatomic congestion games.
\end{exercise}

\begin{exercise}[4]
Investigate the price of anarchy in selfish routing.
\end{exercise}

\begin{exercise}[4]
Prove existence of pure Nash equilibria in finite exact potential games.
\end{exercise}

\begin{exercise}[4]
Study evolutionary stable strategies.
\end{exercise}

\begin{exercise}[4]
Analyze stability of replicator dynamics.
\end{exercise}

%%%%%%%%%%%%%%%%%%%%%%%%%%%%%%%%%%%%%%%%%%%%%%%%%%%%%%%%%%%%
\subsection{Conceptual Summary}
\label{subsec:game-theory-summary}
%%%%%%%%%%%%%%%%%%%%%%%%%%%%%%%%%%%%%%%%%%%%%%%%%%%%%%%%%%%%

Game theory studies strategic interaction among decision makers.

A normal-form game consists of

\[
\boxed{
\text{players}
+
\text{strategy sets}
+
\text{payoff functions}.
}
\]

A Nash equilibrium satisfies

\[
\boxed{
u_i(s_i^*,s_{-i}^*)
\geq
u_i(s_i,s_{-i}^*)
}
\]

for every unilateral deviation \(s_i\).

A mixed strategy is a probability distribution over pure strategies.

For a zero-sum matrix game,

\[
\boxed{
\max_p
\min_q
p^TAq
=
\min_q
\max_p
p^TAq.
}
\]

Sequential games are analyzed using:

\[
\boxed{
\text{extensive form}
+
\text{backward induction}
+
\text{subgame perfection}.
}
\]

Repeated games add future strategic consequences.

Bayesian games add private information and beliefs.

Game theory therefore combines

\[
\boxed{
\text{optimization}
+
\text{probability}
+
\text{economics}
+
\text{strategic reasoning}
+
\text{equilibrium}.
}
\]

%%%%%%%%%%%%%%%%%%%%%%%%%%%%%%%%%%%%%%%%%%%%%%%%%%%%%%%%%%%%
\subsection{Section Connections}
\label{subsec:game-theory-connections}
%%%%%%%%%%%%%%%%%%%%%%%%%%%%%%%%%%%%%%%%%%%%%%%%%%%%%%%%%%%%

\chapterconnection{
Game Theory extends Decision Theory from uncertain but nonstrategic states
of nature to environments containing multiple optimizing agents. Nash
equilibrium formalizes mutual best responses, while zero-sum games connect
directly with linear-programming duality. Extensive-form games connect
with backward induction and Dynamic Programming, Bayesian games combine
strategic behavior with probabilistic beliefs, and congestion games connect
with Network Optimization. The final section of Chapter~10, Optimal
Control, returns to a single decision maker but extends optimization across
continuous-time dynamical systems, leading to variational methods,
Pontryagin's Maximum Principle, and Hamilton--Jacobi--Bellman equations.
}

%%%%%%%%%%%%%%%%%%%%%%%%%%%%%%%%%%%%%%%%%%%%%%%%%%%%%%%%%%%%
% SECTION SOURCE TRACKING
%%%%%%%%%%%%%%%%%%%%%%%%%%%%%%%%%%%%%%%%%%%%%%%%%%%%%%%%%%%%

%------------------------------------------------------------
% 10.15 Game Theory
%------------------------------------------------------------
%
% Source basis:
%   - Standard noncooperative game theory.
%   - Standard zero-sum and extensive-form game theory.
%   - Introductory Bayesian, auction, congestion, and evolutionary games.
%
% Used for:
%   - players, strategies, and payoffs;
%   - normal-form games;
%   - simultaneous and sequential games;
%   - dominant and dominated strategies;
%   - iterated elimination;
%   - best responses;
%   - Nash equilibrium;
%   - Prisoner's Dilemma;
%   - coordination and anti-coordination games;
%   - Battle of the Sexes;
%   - Matching Pennies;
%   - mixed strategies;
%   - expected mixed-strategy payoffs;
%   - indifference principle;
%   - zero-sum games;
%   - maximin, minimax, and saddle points;
%   - von Neumann minimax theorem;
%   - zero-sum linear-programming formulations;
%   - constant-sum and general-sum games;
%   - Pareto efficiency and social welfare;
%   - price of anarchy;
%   - extensive-form games;
%   - information sets;
%   - backward induction;
%   - subgames;
%   - subgame-perfect equilibrium;
%   - credible threats and commitment;
%   - Stackelberg competition preview;
%   - repeated games;
%   - trigger strategies;
%   - Folk Theorem preview;
%   - incomplete information;
%   - Bayesian games;
%   - Bayesian Nash equilibrium;
%   - signaling and screening;
%   - auctions;
%   - first- and second-price sealed-bid auctions;
%   - truthful bidding;
%   - mechanism-design preview;
%   - incentive compatibility;
%   - congestion games;
%   - Wardrop equilibrium preview;
%   - potential games;
%   - evolutionary game theory;
%   - replicator dynamics;
%   - evolutionary stable strategies;
%   - applications and exercises.
%
% Notes:
%   - Optimal Control is developed next in Section 10.16.
%   - Decision Theory appears in Section 10.14.
%   - Network Optimization appears in Section 10.10.
%   - Bayesian probability foundations are developed in Chapters 7 and 8.
%
%%%%%%%%%%%%%%%%%%%%%%%%%%%%%%%%%%%%%%%%%%%%%%%%%%%%%%%%%%%%